% ============================================================================
%  preamble.tex — paketi, lokalizacija, stil
%  Struktura preuzeta iz: AndrijaAdamovic/bachelors-thesis (tex/bachelors_thesis.tex)
%  Oblikovanje usklađeno sa službenim predloškom FSB-a (Predlozak_zavrsni.docx):
%  Times New Roman 12pt, prored 1,5, margine 25mm, numeriranje 1. / 1.1. / 1.1.1.,
%  slike i tablice numerirane kontinuirano, potpisi centrirani bold 11pt.
% ============================================================================

% --- Kodiranje i jezik --------------------------------------------------------
\usepackage[utf8]{inputenc}
\usepackage[T1]{fontenc}
\usepackage[croatian]{babel}
% \usepackage{mathptmx}                   % Times New Roman (po službenom predlošku FSB-a)
\usepackage{lmodern}
\usepackage{setspace}
\onehalfspacing                          % prored 1,5 (stil TEKST)

% --- Stranica ----------------------------------------------------------------
\usepackage[a4paper,margin=25mm,headheight=15pt]{geometry}
% 25mm sa svih strana — po službenom predlošku FSB-a (sectPr: 1418 twips).

% --- Matematika --------------------------------------------------------------
\usepackage{amsmath}
\usepackage{amsthm}
\usepackage{amssymb}
\usepackage{bm}

% --- Tablice -----------------------------------------------------------------
\usepackage{array}
\usepackage{booktabs}
\usepackage{longtable}
\usepackage{multirow}
\usepackage{arydshln}

% --- Slike, kod, ostalo ------------------------------------------------------
\usepackage{graphicx}
\usepackage{caption}
\usepackage{subcaption}
\usepackage{pdfpages}
\usepackage{xcolor}
\usepackage{fancyhdr}
\usepackage{titlesec}
\usepackage{tocloft}
\usepackage{chngcntr}
\usepackage{listings}
\usepackage{algorithm}
\usepackage{algpseudocode}
\usepackage{tikz}
\usetikzlibrary{calc, patterns, decorations.pathmorphing, arrows.meta, positioning}

% --- Mjerne jedinice s hrvatskim decimalnim zarezom ---------------------------
\usepackage{siunitx}
\sisetup{
    output-decimal-marker = {,},
    group-separator       = {.},
    group-minimum-digits  = 4,
    per-mode              = symbol,
    range-units           = single
}
% Upotreba: \SI{0.1}{\milli\metre} -> 0,1 mm ; \SI{3.6}{\kilo\metre\per\hour}

% --- Hiperveze (uvijek zadnje od "običnih" paketa) ---------------------------
\usepackage[
    pdfpagelabels,
    plainpages=false,
    colorlinks=true,
    linkcolor=blue,
    citecolor=blue,
    urlcolor=blue
]{hyperref}
\hypersetup{
    pdfborder={0 0 0},
    linktoc=all,
    pageanchor=false
}

% ============================================================================
%  Hrvatski nazivi fiksnih naslova
% ============================================================================
\addto\captionscroatian{%
    \renewcommand{\contentsname}{SADRŽAJ}%
    \renewcommand{\listfigurename}{POPIS SLIKA}%
    \renewcommand{\listtablename}{POPIS TABLICA}%
    \renewcommand{\bibname}{LITERATURA}%
    \renewcommand{\appendixname}{Prilog}%
    \renewcommand{\figurename}{Slika}%
    \renewcommand{\tablename}{Tablica}%
}

% ============================================================================
%  Tipografija i razmaci
% ============================================================================
\setlength{\parindent}{0pt}
\setlength{\parskip}{3pt}

\setlength{\cftbeforechapskip}{5pt}
\setlength{\cftbeforesecskip}{2pt}
\setlength{\cftbeforesubsecskip}{1pt}
\setlength{\cftbeforefigskip}{5pt}
\setlength{\cftbeforetabskip}{5pt}

\DeclareCaptionLabelFormat{fsb}{#1~#2.}
\captionsetup{
    font          = {small,bf},
    labelformat   = fsb,
    labelsep      = quad,
    justification = centering,
    singlelinecheck = false
}

% --- Numeriranje s točkom na kraju: 1. / 1.1. / 1.1.1. / 1.1.1.1. -------------
\renewcommand{\thechapter}{\arabic{chapter}.}
\renewcommand{\thesection}{\thechapter\arabic{section}.}
\renewcommand{\thesubsection}{\thesection\arabic{subsection}.}
\renewcommand{\thesubsubsection}{\thesubsection\arabic{subsubsection}.}

% --- Slike, tablice i jednadžbe numerirane kontinuirano kroz cijeli rad ------
\counterwithout{figure}{chapter}
\counterwithout{table}{chapter}
\counterwithout{equation}{chapter}

% --- Izgled naslova ---------------------------------------------------------
\titleformat{\chapter}[hang]
{\normalfont\large\bfseries}{\thechapter}{1em}{}
\titleformat{name=\chapter,numberless}[hang]
{\normalfont\large\bfseries}{}{0pt}{}
\titleformat{\section}[hang]
{\normalfont\normalsize\bfseries}{\thesection}{1em}{}
\titleformat{\subsection}[hang]
{\normalfont\normalsize\bfseries\itshape}{\thesubsection}{1em}{}
\titleformat{\subsubsection}[hang]
{\normalfont\normalsize\itshape}{\thesubsubsection}{1em}{}
\titlespacing*{\chapter}{0pt}{0pt}{18pt}
\titlespacing*{\section}{0pt}{14pt}{12pt}
\titlespacing*{\subsection}{0pt}{12pt}{12pt}
\titlespacing*{\subsubsection}{0pt}{12pt}{12pt}
\setcounter{secnumdepth}{3}
\setcounter{tocdepth}{3}

% Šire polje za brojeve u sadržaju (jer brojevi imaju točku na kraju)
\setlength{\cftchapnumwidth}{2.4em}
\setlength{\cftsecnumwidth}{3.2em}
\setlength{\cftsubsecnumwidth}{4.0em}

% ============================================================================
%  Metapodaci rada — IZMIJENI OVDJE
% ============================================================================
\newcommand{\sveuciliste}{Sveučilište u Zagrebu}
\newcommand{\fakultet}{Fakultet strojarstva i brodogradnje}
\newcommand{\vrstarada}{Završni rad}
\newcommand{\autorime}{David Mihić}
\newcommand{\mentorjedan}{Dr. sc. Bojan Šekoranja, mag.ing.}
\newcommand{\mentordva}{}
\newcommand{\mjestogodina}{Zagreb, 2026.}
\newcommand{\naslovrada}{Autonomna interakcija mobilnog manipulatora s okolinom}
\newcommand{\naslovradaen}{Autonomous interaction of a mobile manipulator with the environment}

% ============================================================================
%  Zaglavlje i podnožje
% ============================================================================
\pagestyle{fancy}
\fancyhf{}
\fancyhead[L]{\small \autorime}
\fancyhead[R]{\small \vrstarada}
\fancyfoot[L]{\small\textit{\fakultet{}}}
\fancyfoot[R]{\small\textit{\thepage}}
\renewcommand{\headrulewidth}{0.4pt}
\renewcommand{\footrulewidth}{0.4pt}

\fancypagestyle{plain}{%
    \fancyhf{}
    \fancyhead[L]{\small \autorime}
    \fancyhead[R]{\small \vrstarada}
    \fancyfoot[L]{\small\textit{\fakultet{}}}
    \fancyfoot[R]{\small\textit{\thepage}}
    \renewcommand{\headrulewidth}{0.4pt}
    \renewcommand{\footrulewidth}{0.4pt}
}

% ============================================================================
%  Pomoćne naredbe
% ============================================================================
% Nenumerirano poglavlje koje se ipak pojavi u sadržaju (SAŽETAK, POPIS OZNAKA...)
\newcommand{\frontmatterchapter}[1]{%
    \chapter*{#1}%
    \addcontentsline{toc}{chapter}{#1}%
}

\newcommand{\clearemptypage}{%
    \clearpage
    \thispagestyle{empty}%
}

\newtheorem{definicija}{Definicija}
\newtheorem{teorem}{Teorem}

\newcommand{\norm}[1]{\left\lVert #1 \right\rVert}
\newcommand{\R}{\mathbb{R}}
\newcommand{\transp}{^{\mathsf{T}}}

% Kratice specifične za rad (da se piše konzistentno)
\newcommand{\kmr}{KMR~iiwa}
\newcommand{\lbr}{LBR~iiwa}

% ============================================================================
%  Stil za isječke koda (Python / XML / YAML)
% ============================================================================
\definecolor{codegray}{gray}{0.95}
\definecolor{codekey}{RGB}{0,0,180}
\definecolor{codecom}{RGB}{0,128,0}
\definecolor{codestr}{RGB}{180,0,0}

\lstset{
    basicstyle      = \ttfamily\footnotesize,
    keywordstyle    = \color{codekey}\bfseries,
    commentstyle    = \color{codecom}\itshape,
    stringstyle     = \color{codestr},
    backgroundcolor = \color{codegray},
    numbers         = left,
    numberstyle     = \tiny\color{gray},
    numbersep       = 8pt,
    frame           = single,
    rulecolor       = \color{gray},
    breaklines      = true,
    breakatwhitespace = true,
    showstringspaces  = false,
    tabsize         = 4,
    captionpos      = b,
    inputencoding   = utf8,
    extendedchars   = true,
    literate        = {č}{{\v{c}}}1 {ć}{{\'{c}}}1 {ž}{{\v{z}}}1
        {š}{{\v{s}}}1 {đ}{{d}}1
        {Č}{{\v{C}}}1 {Ć}{{\'{C}}}1 {Ž}{{\v{Z}}}1
        {Š}{{\v{S}}}1 {Đ}{{D}}1
}
\renewcommand{\lstlistingname}{Programski isječak}
